% GENERATED FILE. Edit canonical lessons, then run npm run generate:books.
% canonical-source-hash: fnv1a64:c86fc34262529d50
% canonical-lessons: TE-C362-poradu, TE-C362-tarimikottu, TE-C362-talavancu, TE-C362-sandhi, TE-C362-duta

\chapter{Truce}
\label{ch:te-truce}
\hlchaptermodality{362}

\begin{chapteropening}
\textbf{By the end of this chapter:} \emph{I can say to battle, to drive away, to bow one's head, a truce and an envoy.}

Complete the last lesson of chapter 362: I can say to battle, to drive away, to bow one's head, a truce and an envoy.
\end{chapteropening}

\section[\texorpdfstring{\te{పోరాడు} (pōrāḍu)}{పోరాడు (pōrāḍu)}]{\te{పోరాడు} (pōrāḍu) — to battle; to fight}

\label{lesson:TE-C362-poradu}



\begin{quote}
Before the new one: say the Telugu for to conquer, then the Telugu for a defeat.
\end{quote}

\subsection*{You'll want to know: \te{పోరాడు}}
\textbf{\te{పోరాడు}} — \emph{pōrāḍu} — \textquotedblleft{}to battle; to fight\textquotedblright{}.

\begin{quote}
\te{సైనికులు} \te{చివరి} \te{వరకు} \te{పోరాడారు}.
\end{quote}

\begin{quote}
\emph{sainikulu civari varaku pōrāḍāru.}
\end{quote}

\begin{quote}
The soldiers fought to the end.
\end{quote}

Past (I …): \te{పోరాడాను}, \emph{pōrāḍānu}.

Also to struggle for something: \te{జీవితం} \te{కోసం} \te{పోరాడు}.

\begin{cousinweb}
\textbf{\te{పోరాడు}} (\emph{pōrāḍu}) is native Dravidian: DEDR 4540 lists Telugu \emph{pōru}, \textquotedblleft{}to fight, struggle; a battle\textquotedblright{}, with Tamil \emph{poru}.
\end{cousinweb}

\begin{grammarlens}[title={What you've built}]
One more word for telling what happened.
\end{grammarlens}

\subsection*{Your turn}
\begin{itemize}
\raggedright
  \item \emph{Say it:} \emph{pōrāḍu}
  \item \emph{Say it:} \emph{pōrāḍu}, once more
  \item \emph{Say it:} \emph{sainikulu civari varaku pōrāḍāru.} — The soldiers fought to the end.
  \item \emph{Recall:} say the Telugu for to conquer, then the Telugu for a defeat, then say \emph{pōrāḍu} again
\end{itemize}

\begin{tcolorbox}[breakable,colback=teal!4,colframe=teal!35!black,title={Before you move on}]
What does \te{పోరాడు} mean? (\textquotedblleft{}To battle; to fight\textquotedblright{}.) Say it once more.
\end{tcolorbox}
\section[\texorpdfstring{\te{తరిమికొట్టు} (tarimikoṭṭu)}{తరిమికొట్టు (tarimikoṭṭu)}]{\te{తరిమికొట్టు} (tarimikoṭṭu) — to drive away; to chase off}

\label{lesson:TE-C362-tarimikottu}



\begin{quote}
Before the new one: say the Telugu for a defeat, then the Telugu for to battle.
\end{quote}

\subsection*{You'll want to know: \te{తరిమికొట్టు}}
\textbf{\te{తరిమికొట్టు}} — \emph{tarimikoṭṭu} — \textquotedblleft{}to drive away; to chase off\textquotedblright{}.

\begin{quote}
\te{గ్రామస్తులు} \te{దొంగలను} \te{తరిమికొట్టారు}.
\end{quote}

\begin{quote}
\emph{grāmastulu doṅgalanu tarimikoṭṭāru.}
\end{quote}

\begin{quote}
The villagers drove the robbers away.
\end{quote}

Past (I …): \te{తరిమికొట్టాను}, \emph{tarimikoṭṭānu}.

\te{తరుము}, to chase, with \te{కొట్టు}: to chase right off.

\begin{cousinweb}
\textbf{\te{తరిమికొట్టు}} (\emph{tarimikoṭṭu}) is a Telugu compound: \emph{tarumu}, \textquotedblleft{}to chase\textquotedblright{}, and native \emph{koṭṭu}, \textquotedblleft{}to beat, strike\textquotedblright{} (DEDR 2063).
\end{cousinweb}

\begin{grammarlens}[title={What you've built}]
One more word for telling what happened.
\end{grammarlens}

\subsection*{Your turn}
\begin{itemize}
\raggedright
  \item \emph{Say it:} \emph{tarimikoṭṭu}
  \item \emph{Say it:} \emph{tarimikoṭṭu}, once more
  \item \emph{Say it:} \emph{grāmastulu doṅgalanu tarimikoṭṭāru.} — The villagers drove the robbers away.
  \item \emph{Recall:} say the Telugu for a defeat, then the Telugu for to battle, then say \emph{tarimikoṭṭu} again
\end{itemize}

\begin{tcolorbox}[breakable,colback=teal!4,colframe=teal!35!black,title={Before you move on}]
What does \te{తరిమికొట్టు} mean? (\textquotedblleft{}To drive away; to chase off\textquotedblright{}.) Say it once more.
\end{tcolorbox}
\section[\texorpdfstring{\te{తలవంచు} (talavañcu)}{తలవంచు (talavañcu)}]{\te{తలవంచు} (talavañcu) — to bow one's head; to submit}

\label{lesson:TE-C362-talavancu}



\begin{quote}
Before the new one: say the Telugu for to battle, then the Telugu for to drive away.
\end{quote}

\subsection*{You'll want to know: \te{తలవంచు}}
\textbf{\te{తలవంచు}} — \emph{talavañcu} — \textquotedblleft{}to bow one's head; to submit\textquotedblright{}.

\begin{quote}
\te{చివరికి} \te{రాజు} \te{శత్రువుకు} \te{తలవంచాడు}.
\end{quote}

\begin{quote}
\emph{civariki rāju śatruvuku talavañcāḍu.}
\end{quote}

\begin{quote}
In the end the king submitted to the enemy.
\end{quote}

Past (I …): \te{తలవంచాను}, \emph{talavañcānu}.

Literally to bend the head: \te{తల}, head, and \te{వంచు}, to bend.

\begin{cousinweb}
\textbf{\te{తలవంచు}} (\emph{talavañcu}) is a native compound: \emph{tala}, \textquotedblleft{}the head\textquotedblright{} (DEDR 3103), and \emph{vañcu}, \textquotedblleft{}to bend, humble\textquotedblright{} (DEDR 5335).
\end{cousinweb}

\begin{grammarlens}[title={What you've built}]
One more word for telling what happened.
\end{grammarlens}

\subsection*{Your turn}
\begin{itemize}
\raggedright
  \item \emph{Say it:} \emph{talavañcu}
  \item \emph{Say it:} \emph{talavañcu}, once more
  \item \emph{Say it:} \emph{civariki rāju śatruvuku talavañcāḍu.} — In the end the king submitted to the enemy.
  \item \emph{Recall:} say the Telugu for to battle, then the Telugu for to drive away, then say \emph{talavañcu} again
\end{itemize}

\begin{tcolorbox}[breakable,colback=teal!4,colframe=teal!35!black,title={Before you move on}]
What does \te{తలవంచు} mean? (\textquotedblleft{}To bow one's head; to submit\textquotedblright{}.) Say it once more.
\end{tcolorbox}
\section[\texorpdfstring{\te{సంధి} (sandhi)}{సంధి (sandhi)}]{\te{సంధి} (sandhi) — a truce; a treaty}

\label{lesson:TE-C362-sandhi}



\begin{quote}
Before the new one: say the Telugu for to drive away, then the Telugu for to bow one's head.
\end{quote}

\subsection*{You'll want to know: \te{సంధి}}
\textbf{\te{సంధి}} — \emph{sandhi} — \textquotedblleft{}a truce; a treaty\textquotedblright{}.

\begin{quote}
\te{రెండు} \te{రాజ్యాలు} \te{సంధి} \te{చేసుకున్నాయి}.
\end{quote}

\begin{quote}
\emph{reṇḍu rājyālu sandhi cēsukunnāyi.}
\end{quote}

\begin{quote}
The two kingdoms made a truce.
\end{quote}

\te{సంధి} \te{చేసుకో}, to make peace; in grammar, \te{సంధి} is the joining of sounds.

\begin{cousinweb}
\textbf{\te{సంధి}} (\emph{sandhi}) is borrowed from Sanskrit \emph{saṁdhi}, \textquotedblleft{}a joint, a point of union; a compact\textquotedblright{} (CDIAL 12913).
\end{cousinweb}

\begin{grammarlens}[title={What you've built}]
One more word for telling what happened.
\end{grammarlens}

\subsection*{Your turn}
\begin{itemize}
\raggedright
  \item \emph{Say it:} \emph{sandhi}
  \item \emph{Say it:} \emph{sandhi}, once more
  \item \emph{Say it:} \emph{reṇḍu rājyālu sandhi cēsukunnāyi.} — The two kingdoms made a truce.
  \item \emph{Recall:} say the Telugu for to drive away, then the Telugu for to bow one's head, then say \emph{sandhi} again
\end{itemize}

\begin{tcolorbox}[breakable,colback=teal!4,colframe=teal!35!black,title={Before you move on}]
What does \te{సంధి} mean? (\textquotedblleft{}A truce; a treaty\textquotedblright{}.) Say it once more.
\end{tcolorbox}
\section[\texorpdfstring{\te{దూత} (dūta)}{దూత (dūta)}]{\te{దూత} (dūta) — an envoy; a messenger}

\label{lesson:TE-C362-duta}



\begin{quote}
Before the new one: say the Telugu for to bow one's head, then the Telugu for a truce.
\end{quote}

\subsection*{You'll want to know: \te{దూత}}
\textbf{\te{దూత}} — \emph{dūta} — \textquotedblleft{}an envoy; a messenger\textquotedblright{}.

\begin{quote}
\te{రాజు} \te{శత్రువు} \te{దగ్గరకు} \te{దూతను} \te{పంపాడు}.
\end{quote}

\begin{quote}
\emph{rāju śatruvu daggaraku dūtanu pampāḍu.}
\end{quote}

\begin{quote}
The king sent an envoy to the enemy.
\end{quote}

Plural: \te{దూతలు}, \emph{dūtalu}.

\begin{cousinweb}
\textbf{\te{దూత}} (\emph{dūta}) is borrowed from Sanskrit \emph{dūta}, \textquotedblleft{}a messenger\textquotedblright{} (CDIAL 6488).
\end{cousinweb}

\begin{grammarlens}[title={What you've built}]
One more word for telling what happened.
\end{grammarlens}

\subsection*{Your turn}
\begin{itemize}
\raggedright
  \item \emph{Say it:} \emph{dūta}
  \item \emph{Say it:} \emph{dūta}, once more
  \item \emph{Say it:} \emph{rāju śatruvu daggaraku dūtanu pampāḍu.} — The king sent an envoy to the enemy.
  \item \emph{Recall:} say the Telugu for to bow one's head, then the Telugu for a truce, then say \emph{dūta} again
\end{itemize}

\begin{tcolorbox}[breakable,colback=teal!4,colframe=teal!35!black,title={Before you move on}]
What does \te{దూత} mean? (\textquotedblleft{}An envoy; a messenger\textquotedblright{}.) Say it once more.
\end{tcolorbox}
